\uuid{lUM8}
\titre{DL usuels/Etude locale}
\auteur{hanani ; mzari}
\organisation{exo7}
\licence{CC BY-NC-SA 4.0}
\datecreate{2020-01-20}
\niveau{L1}
\module{Analyse}
\chapitre{Développement limité}
\sousChapitre{Applications}
\difficulte{2}
\theme{Applications des DL, DL usuels, Etude locale}
\exo7qcmid{lille-2.developpements-limites.033}

\contenu{
\texte{
Soit $a\in \Rr$ et $f$ une fonction. On dit que $f(x)$ est un $o\left((x-a)^n\right)$ au voisinage de $a$ et on écrit $f(x)=o((x-a)^n)$ si $\displaystyle \lim _{x\to a}\frac{f(x)}{(x-a)^n}=0$. On note par $DL_n(a)f(x)$ le développement limité de $f$ à l'ordre $n$ en $a$. On note aussi $DL_n(a^+)f(x)$ (resp. $DL_n(a^-)f(x)$) le développement limité de $f$ à l'ordre $n$ à droite en $a$ (resp. à gauche en $a$).
}

\begin{question}
\qclassification{125.03}
\qtags{motcle=DL usuels/Etude locale}
Soit $f$ telle que $\displaystyle f(x)=\frac{\mathrm{e}^{x}-1}{x+x^2}$ si $x\neq 0$ et $f(0)=1$. On note $T_0$ la tangente au graphe de $f$ au point d'abscisse $0$ lorsqu'elle existe. Parmi les affirmations suivantes, cocher celles qui sont vraies :
\begin{answers}
\bad{$\displaystyle DL_2(0)f(x)=1-\frac{x}{2}+\frac{x^2}{2}+o(x^2)$.}
\good{$\displaystyle \lim _{x\to 0}f(x)=1$ et $f$ est continue en $0$.}
\good{$f$ est dérivable en $0$ et $\displaystyle f'(0)=-\frac{1}{2}$.}
\bad{$T_0$ est la droite d'équation $\displaystyle y=1-\frac{x}{2}$ et le graphe de $f$ est en dessous de $T_0$ au voisinage de $0$.}
\end{answers}
\begin{explanations}
Pour écrire le $DL_2(0)f(x)$, on utilise le $DL_3(0)(\mathrm{e}^{x}-1)$. La division suivant les puissances croissantes donne :
$$\displaystyle DL_2(0)f(x)=1-\frac{x}{2}+\frac{2x^2}{3}+o(x^2).$$
Ainsi $\displaystyle \lim _{x\to 0}f(x)=1$ et $f$ est continue en $0$ et, puisque $f$ admet un $DL_1(0)$, $f$ est dérivable en $0$. De plus, $T_0$ est la droite d'équation $\displaystyle y=1-\frac{x}{2}$ et puisque $\displaystyle f(x)-y=\frac{2x^2}{3}+o(x^2)\geq 0$ au voisinage de $0$ : le graphe de $f$ est au dessus de $T_0$ au voisinage de $0$.
\end{explanations}
\end{question}
}
